\section{Математическое описание}

\subsection{Граф и его определение}

Граф~--- пара $G = (V, E)$, где $V$~--- конечное непустое множество вершин,
$E \subseteq V \times V$~--- множество рёбер. Граф называется ориентированным,
если каждое ребро~--- упорядоченная пара, и неориентированным в противном
случае. Ациклический граф не содержит циклов; связный ациклический
неориентированный граф является деревом, ориентированный~--- DAG
(directed acyclic graph).

Взвешенный граф задаёт функцию весов $w : E \to \mathbb{R}$, ставящую
каждому ребру вещественное число. В данной работе веса генерируются
по распределению Вейбулла.

\subsection{Распределение классического Вейбулла}

Согласно справочнику Вадзинского, распределение Вейбулла задаётся
плотностью вероятности:
\[
f(x) = \frac{c}{a} \left(\frac{x}{a}\right)^{c-1}
        \exp\!\left(-\left(\frac{x}{a}\right)^c\right), \quad x > 0,
\]
где $a > 0$~--- параметр масштаба (\textit{характерное время жизни}),
$c > 0$~--- параметр формы. При любом $c > 0$
$P\{X < a\} = 1 - e^{-1} \approx 0.6321$.

Для генерации случайных значений используется метод обратного преобразования.
Сначала генерируется равномерное число $U \sim \mathrm{Uniform}(0,1)$,
после чего применяется обратная функция распределения Вейбулла:
\[
X = a\,(-\ln(1 - U))^{1/c}.
\]
Это позволяет получить значения, распределённые по закону Вейбулла, что
делает генерацию параметров графа более реалистичной и менее равномерной
по сравнению с обычным генератором случайных чисел.

В реализации используются два независимых экземпляра распределения
с параметрами \texttt{a = 10, c = 2}: один для генерации структуры графа
(степеней вершин), другой~--- для генерации весов рёбер.

\subsection{Генерация связного ациклического графа по степеням вершин}

Пользователь задаёт число вершин~$n$, тип ориентированности и знак весов.
Число рёбер не вводится напрямую.

\textbf{Шаг~1: генерация целевых степеней.}
Для каждой вершины $i \in \{0,\ldots,n-1\}$ независимо сэмплируется
значение по распределению Вейбулла, округляемое до целого и зажатое
в диапазон $[1,\,\lfloor\log n\rfloor]$:
\[
    d_i = \max\!\bigl(1,\; \min\!\bigl(\lfloor X_i \rfloor,\; \lfloor\log n\rfloor\bigr)\bigr),
    \quad X_i \sim \mathrm{Weibull}(a, c).
\]

\textbf{Шаг~2: формирование остовного дерева.}
Задаётся перестановка вершин $\pi$ (в данной реализации --- тождественная,
чтобы сохранялась совместимость с топологической сортировкой).
Для каждой вершины $\pi[i]$ при $i \geq 1$ выбирается предок случайно
из $\{\pi[0], \ldots, \pi[i-1]\}$ и добавляется ребро.
Построение по индексам $j < i$ гарантирует ациклическость: ни одно ребро
не идёт «назад» в топологическом порядке.

\textbf{Шаг~3: добирание дополнительных рёбер.}
Для каждой вершины $\pi[i]$ с оставшимся бюджетом степени $d_i > 0$
формируется множество допустимых целей $\{j \mid j > i,\; (\pi[i],\pi[j]) \notin E\}$.
Это множество перемешивается с помощью алгоритма Фишера-Йетса, где индексы
выбираются по распределению Вейбулла (\textit{weibullShuffle}), после чего
добавляются первые $\min(d_i, |\text{допустимых}|)$ рёбер.

Ориентированность при добавлении рёбер учитывается явно: для неориентированных
графов каждое ребро добавляется симметрично.

\subsection{Эксцентриситет, центр и диаметральные вершины}

\textit{Эксцентриситетом} $e(v)$ вершины $v$ в связном графе $G(V, E)$
называется максимальное расстояние от вершины $v$ до других вершин графа $G$:
\[
e(v) \stackrel{\text{Def}}{=} \max_{u \in V} d(v, u).
\]

\textit{Радиусом} $R(G)$ графа $G$ называется наименьший из эксцентриситетов вершин:
\[
R(G) \stackrel{\text{Def}}{=} \min_{v \in V} e(v).
\]

Вершина $v$ называется \textit{центральной}, если её эксцентриситет совпадает
с радиусом графа, $e(v) = R(G)$. Множество центральных вершин называется
\textit{центром} графа:
\[
C(G) \stackrel{\text{Def}}{=} \{v \in V \mid e(v) = R(G)\}.
\]

Диаметральные вершины~--- с максимальным эксцентриситетом:
\[
D(G) = \{v \mid e(v) = \mathrm{diam}(G)\}.
\]

Расстояния вычисляются алгоритмом обхода в ширину ($BFS$), запускаемым
из каждой вершины; алгоритмическая сложность линейная ($O(V + E)$) для
одного прохода, а итоговая сложность квадратична $O(V(V + E))$.

\subsection{Метод Шимбелла}

\begin{enumerate}
    \item Операция «умножения» двух величин $a$ и $b$ при возведении матрицы
    в степень соответствует их алгебраической сумме:
    \begin{equation}
        a \cdot b = b \cdot a \Rightarrow a + b = b + a, \qquad
        a \cdot 0 = 0 \cdot a = 0 \Rightarrow a + 0 = 0.
    \end{equation}

    \item Операция «сложения» двух величин $a$ и $b$ заменяется выбором
    минимального (максимального) элемента из ненулевых:
    \begin{equation}
        a + b = b + a = \min(\max)\{a, b\}.
    \end{equation}
    Нули при этом игнорируются: минимальный или максимальный элемент
    выбирается из ненулевых значений. Нуль в результате операции может быть
    получен лишь тогда, когда все слагаемые --- нулевые.
\end{enumerate}

Элемент матрицы $A^{(2)}$, полученной возведением $A$ в степень~2,
вычисляется как:
\[
\omega_{ij}^{(2)} = \min(\max) \left\{
    \omega_{i1}^{(1)} + \omega_{1j}^{(1)},\;
    \omega_{i2}^{(1)} + \omega_{2j}^{(1)},\;
    \ldots,\;
    \omega_{in}^{(1)} + \omega_{nj}^{(1)}
\right\}.
\]
Матрица $A^{(k)}$, вычисленная последовательным умножением, содержит длины
путей из ровно $k$ рёбер для всех пар вершин. Временная сложность: $O(k \cdot n^3)$,
пространственная: $O(n^2)$.

\subsection{Упорядочивание графа по алгоритму Фалкерсона}

Алгоритм Фалкерсона --- графический метод топологической сортировки DAG,
состоящий из следующих шагов.

\begin{enumerate}
    \item Найти вершины графа, в которые не входит ни одна дуга.
    Они образуют первую группу. Пронумеровать вершины группы в произвольном порядке.

    \item Вычеркнуть все пронумерованные вершины и дуги, из них исходящие.
    В получившемся графе найдётся, по крайней мере, одна вершина, в которую
    не входит ни одна дуга. Этой вершине, входящей во вторую группу, присвоить
    очередной номер и так далее. Второй шаг повторять до тех пор, пока не будут
    упорядочены все вершины.
\end{enumerate}

Если на каком-либо шаге вершин без входящих дуг не оказалось, граф содержит
цикл и алгоритм неприменим. Алгоритмическая сложность: $O(V^2)$ в матричном
представлении, поскольку на каждом шаге выполняется просмотр всей матрицы
смежности для поиска вершин без входящих дуг.

\subsection{Нахождение кратчайшего пути. Алгоритм Флойда-Уоршалла}

Алгоритм Флойда-Уоршалла находит кратчайшие пути между \textit{всеми} парами
вершин в ориентированном взвешенном графе.

Для хранения информации о путях используется матрица $H$ размера $n \times n$:
\[
H[i, j] =
\begin{cases}
    k, & \text{если $k$~--- первая вершина, достигаемая на кратчайшем
    пути из $i$ в $j$;} \\
    -1, & \text{если из вершины $i$ в вершину $j$ нет пути.}
\end{cases}
\]

Основной цикл: для каждой промежуточной вершины $k$ проверяется
\[
    \text{если } d[i][k] + d[k][j] < d[i][j], \quad
    \text{то } d[i][j] \leftarrow d[i][k] + d[k][j], \quad H[i][j] \leftarrow H[i][k].
\]

\textbf{\textit{Если в $G$ есть цикл с отрицательным весом, то решения
поставленной задачи не существует.}}

Алгоритмическая сложность: $O(n^3)$, пространственная: $O(n^2)$.

\subsection{Определение сети}

Пусть $G = (V, E)$~--- орграф с одним истоком $s \in V$ и одним стоком
$t \in V$, где $s \neq t$. \textbf{Сетью} $S = (G;\, c)$ называется орграф
$G$ с заданной функцией $c : E \to \mathbb{R}_+$, сопоставляющей каждой
дуге $e \in E$ неотрицательное действительное число $c(e)$~--- пропускную
способность.

\subsection{Поток в сети}

\textbf{Потоком} в сети $S$ называется функция $f : E \to \mathbb{R}_+$,
для которой выполняются два условия.

Первое~--- ограничение пропускной способности: для каждой дуги $e \in E$
\[
    f(e) \leqslant c(e).
\]

Второе~--- закон сохранения потока: для каждой вершины $v \in V$,
$v \neq s, t$,
\[
    D_f(v) = \sum_{e \in A(v)} f(e) - \sum_{e \in B(v)} f(e) = 0,
\]
где $A(v)$~--- множество дуг, исходящих из $v$,
$B(v)$~--- множество дуг, входящих в $v$.

\textbf{Величиной потока} $f$ называется число
\[
    p_f = \sum_{e \in A(s)} f(e).
\]

\subsection{Алгоритм Форда--Фалкерсона}

\textbf{Теорема} (Л.\,Р.\,Форд, Д.\,Р.\,Фалкерсон, 1962).
\textit{Величина максимального потока $p_{\max}$ в сети $S$ совпадает
с минимальной пропускной способностью $r_{\min}$ её разрезов.}

Алгоритм итеративно находит увеличивающий путь от истока к стоку
в остаточной сети. Величина дополнительного потока по найденному пути:
\[
    \Delta f = \min_{(u,v) \in P} r(u,v).
\]
Для каждого ребра $(u,v)$ пути $P$ поток обновляется:
\[
    f(u,v) \leftarrow f(u,v) + \Delta f, \qquad
    f(v,u) \leftarrow f(v,u) - \Delta f.
\]
Алгоритм завершается, когда увеличивающий путь не существует.
В реализации поиск пути выполняется через BFS (алгоритм Эдмондса--Карпа),
что гарантирует сложность $O(V \cdot E^2)$.

\subsection{Нахождение заданного потока минимальной стоимости}

Рассмотрим задачу определения потока заданной величины $\theta$ от $s$ к
$t$ в сети $G = (S, U, \Omega, D)$, в которой каждая дуга $(x_i, x_j)$
характеризуется пропускной способностью $c(x_i, x_j) \in \Omega$ и
неотрицательной стоимостью $d(x_i, x_j) \in D$.

Математическая модель задачи: минимизировать
\[
    \sum_{(x_i,\, x_j) \in U} d(x_i, x_j) \cdot \varphi(x_i, x_j),
\]
при ограничениях:

\begin{enumerate}
    \item $0 \leqslant \varphi(x_i, x_j) \leqslant c(x_i, x_j)$
          для всех $(x_i, x_j) \in U$;
    \item для истока $s$:
          $\displaystyle\sum_{x_j} \varphi(s, x_j) - \sum_{x_j} \varphi(x_j, s) = \theta$
    \item для стока $t$:
          $\displaystyle\sum_{x_j} \varphi(t, x_j) - \sum_{x_j} \varphi(x_j, t) = -\theta$
    \item для любой другой вершины: условие сохранения потока равно нулю.
\end{enumerate}

\subsection{Теорема Кирхгофа}

Пусть $G = (V, E)$~--- связный неориентированный граф без петель и кратных
рёбер, $|V| = n$. \textit{Матрица Кирхгофа} $B = D - A(G)$,
где $D = \mathrm{diag}(\deg v_1, \ldots, \deg v_n)$~--- диагональная матрица
степеней вершин, $A(G)$~--- матрица смежности:
\[
    b_{ij} =
    \begin{cases}
        \deg v_i, & i = j, \\
        -1,       & i \neq j,\ v_i \sim v_j, \\
        0,        & i \neq j,\ v_i \not\sim v_j.
    \end{cases}
\]

\textbf{Теорема} (матричная теорема Кирхгофа). \textit{Число различных
остовных деревьев связного графа $G$ равно определителю любого минора
матрицы Кирхгофа:}
\[
    \tau(G) = \det B_{ii},
\]
\textit{где $B_{ii}$~--- матрица $B$ с удалённой $i$-й строкой и $i$-м столбцом.}

В реализации определитель вычисляется методом Гаусса с целочисленной
арифметикой~--- ошибок округления нет.
Алгоритмическая сложность: $O(n^3)$, пространственная: $O(n^2)$.

\subsection{Алгоритм Краскала}

Будем последовательно строить подграф $F$ графа $G$ («растущий лес»),
пытаясь на каждом шаге достроить $F$ до некоторого MST. Начнём с того, что
включим в $F$ все вершины графа $G$. Теперь будем обходить множество $E(G)$
в порядке неубывания весов рёбер. Если очередное ребро $e$ соединяет вершины
одной компоненты связности $F$, то добавление его в остов приведёт к возникновению
цикла --- $e$ не включается. Иначе $e$ соединяет разные компоненты связности,
и из леммы о безопасном ребре следует, что $e$ является безопасным~--- добавляем
его в $F$.

Для проверки цикличности в реализации используется BFS по текущему остову.
Алгоритмическая сложность: $O(E^2)$ (проверка через BFS), пространственная: $O(V + E)$.

\subsection{Код Прюфера}

\textbf{Кодирование.} Пусть $T = (V, E_T)$~--- дерево с вершинами
$V = \{0, \ldots, n-1\}$. Код Прюфера~--- последовательность
$A = (A[1], \ldots, A[n-2])$: на каждом из $n - 2$ шагов выбирается
висячая вершина $v$ с минимальным номером ($\deg v = 1$), в код записывается
её единственный сосед $\Gamma(v)$, после чего вершина $v$ удаляется из дерева.
Вместе с кодом сохраняются веса рёбер.

\textbf{Декодирование.} По коду $A = (A[1], \ldots, A[n-2])$
восстанавливается дерево $T = (V, E_T)$. Поддерживается множество
«свободных» вершин $B$. На каждом шаге~$i$ выбирается наименьшая вершина
$v \in B$, не встречающаяся в остатке кода $A[i..n-2]$; добавляется ребро
$(v, A[i])$, вершина $v$ удаляется из $B$. Веса восстанавливаются из
сохранённого словаря. Сложность кодирования и декодирования: $O(n^2)$,
пространственная: $O(n)$.

\subsection{Максимальное независимое множество вершин}

\textit{Независимым множеством} графа $G = (V, E)$ называется подмножество
$S \subseteq V$ такое, что никакие два его элемента не соединены ребром:
\[
    \forall u, v \in S : (u, v) \notin E.
\]
\textit{Максимальным независимым множеством} (MIS) называется такое
независимое множество $S^*$, что $|S^*| \geq |S|$ для любого другого
независимого множества $S$.

Нахождение MIS является NP-трудной задачей в общем случае.
В данной работе реализован точный алгоритм на основе полного перебора
с отсечениями (backtracking): рекурсивно перебираются кандидаты,
каждый из которых добавляется в текущее множество $S$, если он
не смежен ни с одной уже включённой вершиной. После добавления
вершины $v$ из кандидатов исключаются все её соседи.
Сложность: $O(2^n)$ в худшем случае, пространственная: $O(n)$.
